\documentclass[a4paper,11pt]{article}
% ---------- Кириллица: работает и в pdfLaTeX, и в XeLaTeX/LuaLaTeX ----------
\usepackage{iftex}
\ifPDFTeX
  \usepackage[T2A]{fontenc}
  \usepackage[utf8]{inputenc}
  \usepackage[russian]{babel}
\else
  \usepackage{fontspec}
  \IfFontExistsTF{CMU Serif}{\setmainfont{CMU Serif}}{%
    \IfFontExistsTF{Times New Roman}{\setmainfont{Times New Roman}}{\setmainfont{FreeSerif}}}
  \IfFontExistsTF{CMU Typewriter Text}{%
    \setmonofont{CMU Typewriter Text}\newfontfamily\cyrillicfonttt{CMU Typewriter Text}}{%
    \IfFontExistsTF{Consolas}{%
      \setmonofont{Consolas}\newfontfamily\cyrillicfonttt{Consolas}}{%
      \setmonofont{FreeMono}\newfontfamily\cyrillicfonttt{FreeMono}}}
  \usepackage{polyglossia}
  \setdefaultlanguage{russian}
\fi
\usepackage{amsmath,amssymb,mathtools}
\usepackage[a4paper,margin=2cm]{geometry}
\usepackage[dvipsnames]{xcolor}
\usepackage{tikz}
\usetikzlibrary{arrows.meta,calc,patterns,decorations.pathreplacing,babel,
  shapes.geometric,shapes.misc,shapes.symbols,positioning}
\usepackage[most]{tcolorbox}
\usepackage{enumitem}
\usepackage{titlesec}
\usepackage{array}
\usepackage{fancyvrb}
\usepackage[unicode,hidelinks]{hyperref}

\emergencystretch=3em
\tolerance=1500
\setlength{\parindent}{0pt}
\setlength{\parskip}{4pt}
\setlist{nosep,leftmargin=*,topsep=2pt}

% ---------- оформление (как в angem.tex) ----------
\definecolor{secred}{RGB}{190,40,45}
\definecolor{subteal}{RGB}{0,110,120}
\definecolor{leadcol}{RGB}{0,120,110}
\titleformat{\section}{\Large\bfseries\color{secred}}{\thesection.}{0.5em}{}
\titleformat{\subsection}{\large\bfseries\color{subteal}}{\thesubsection.}{0.5em}{}
\titlespacing*{\section}{0pt}{14pt}{6pt}
\titlespacing*{\subsection}{0pt}{10pt}{4pt}

\newtcolorbox{opr}{breakable,enhanced,colback=blue!4,colframe=blue!55!black,
  boxrule=0pt,leftrule=2.2pt,arc=0pt,left=6pt,right=4pt,top=3pt,bottom=3pt,
  before skip=6pt,after skip=6pt}
\newtcolorbox{teor}[1]{breakable,enhanced,colback=orange!7,colframe=orange!70!black,
  boxrule=0pt,leftrule=2.2pt,arc=0pt,left=6pt,right=4pt,top=3pt,bottom=3pt,
  before skip=6pt,after skip=6pt,before upper={\textbf{Теорема #1.}\ }}
\newtcolorbox{hwbox}{breakable,enhanced,colback=gray!6,colframe=gray!60,
  boxrule=0.4pt,arc=1pt,left=6pt,right=6pt,top=3pt,bottom=3pt,fontupper=\small}

\newcommand{\lead}[1]{\textbf{\color{leadcol}#1}\ }
\newcommand{\q}{\textbf{\color{red!70!black}?}\ }
\newcommand{\imp}{\textbf{\color{red!70!black}!}\ }
\newcommand{\zam}{\textbf{Замечание.}\ }
\newcommand{\prim}{\textbf{Пример.}\ }
\newcommand{\primn}[1]{\textbf{Пример #1.}\ }
\newcommand{\ispr}[1]{\footnote{\textit{Исправлено.} #1}}
\renewcommand{\thefootnote}{\textcolor{red!75!black}{\arabic{footnote}}}
\newcommand{\hwlist}[1]{{\small\color{gray}[#1]}}
\newcommand{\qedr}{\hfill$\blacksquare$}
\DeclareRobustCommand{\cn}[1]{\tikz[baseline=(c.base)]\node[draw,circle,inner sep=0.7pt,minimum size=1.15em,font=\footnotesize\bfseries] (c) {#1};}

% Подпись и рисунок под ней не разрываются между страницами
\newenvironment{schema}{\par\noindent\begin{minipage}{\linewidth}\setlength{\parskip}{4pt}}%
  {\end{minipage}\par\smallskip}

% ---------- код ----------
% Код пишется в окружении code (как есть, без экранирования):
%   \begin{code}
%   print('Привет')
%   \end{code}
% Короткий код в тексте --- \texttt{...} (символы _ % # { } экранировать).
\definecolor{codebg}{RGB}{245,247,248}
\newtcolorbox{codebox}{breakable,enhanced,colback=codebg,colframe=subteal!75,
  boxrule=0pt,leftrule=1.5pt,arc=0pt,left=6pt,right=4pt,top=1pt,bottom=1pt,
  before skip=5pt,after skip=5pt}
\newenvironment{code}%
  {\VerbatimEnvironment\begin{codebox}\begin{Verbatim}[fontsize=\small]}%
  {\end{Verbatim}\end{codebox}}

% ---------- блоки схем алгоритмов (ГОСТ 19.701-90) ----------
% 1 клетка = 0,5 см: блок 6 x 2 клетки = 3 x 1 см
\tikzset{
  sa/.style={line width=0.6pt,font=\small,node distance=5mm,>={Stealth[length=2mm]}},
  blk/.style={draw,align=center,minimum width=3cm,minimum height=1cm,inner xsep=2pt,inner ysep=3.5pt},
  proc/.style={blk,rectangle},                                   % процесс
  term/.style={blk,rounded rectangle},                           % начало, конец
  io/.style={blk,trapezium,trapezium left angle=70,trapezium right angle=110,
    trapezium stretches body},                                   % ввод, вывод
  dec/.style={blk,diamond,aspect=2.4,minimum width=3cm,minimum height=1.2cm,
    inner sep=1pt},                                              % решение
  prep/.style={blk,signal,signal to=east and west,signal pointer angle=120},
                                                                 % цикл (старый ГОСТ)
  lpbeg/.style={blk,chamfered rectangle,chamfered rectangle corners={north west,north east},
    chamfered rectangle sep=2.5mm},                              % начало цикла (новый ГОСТ)
  lpend/.style={blk,chamfered rectangle,chamfered rectangle corners={south west,south east},
    chamfered rectangle sep=2.5mm},                              % конец цикла (новый ГОСТ)
  pred/.style={proc,path picture={\draw
    ([xshift=2.5mm]path picture bounding box.north west)--([xshift=2.5mm]path picture bounding box.south west)
    ([xshift=-2.5mm]path picture bounding box.north east)--([xshift=-2.5mm]path picture bounding box.south east);}},
                                                                 % вызов подпрограммы
  yes/.style={font=\footnotesize,inner sep=1pt},                 % подписи «да», «нет»
  cmt/.style={font=\small,align=left},                           % комментарий
}
\newcommand{\dop}[1]{\footnote{\textit{Дописано.} #1}}
% Цикл по старому ГОСТу: линия возврата от последнего блока тела (#3) в левую вершину
% шестиугольника (#2), выход из правой вершины. Точка выхода --- координата (#2-out).
% Необязательный параметр --- отступ линий от блоков, см.
\newcommand{\cycleold}[3][0.45]{%
  \draw[->] (#3.south) -- ++(0,-0.3) -| ([xshift=-#1 cm]#2.west) -- (#2.west);
  \draw (#2.east) -- ++(#1,0) |- ([yshift=-0.6cm]#3.south)
    node[inner sep=0pt,minimum size=0pt] (#2-out) {};}
% комментарий по ГОСТ: пунктир от точки (#1) к скобке с текстом #2
\newcommand{\gostcomment}[3][1.2]{%
  \draw[dashed] (#2) -- ++(#1,0) coordinate (cm@a);
  \draw ([shift={(0.18,0.32)}]cm@a) -- ([shift={(0,0.32)}]cm@a) -- ([shift={(0,-0.32)}]cm@a) -- ([shift={(0.18,-0.32)}]cm@a);
  \node[cmt,anchor=west] at ([xshift=0.15cm]cm@a) {#3};}


\begin{document}
\begin{center}
{\LARGE\bfseries\color{secred}Программирование}\\[3pt]
{\large Конспект лекций}
\end{center}
\vspace{-2mm}

\tableofcontents
\newpage


% =====================================================================
\section{Арифметические выражения}
% =====================================================================
\emph{Операция}: \texttt{17 + 2}. \ \emph{Оператор} (команда, инструкция).

\lead{Арифметическое выражение может состоять из:}
\begin{enumerate}[label=\arabic*)]
\item \textbf{операций}:\\[2pt]
\begin{tabular}{@{\qquad}>{\ttfamily}l@{\quad}l@{\qquad\qquad}>{\ttfamily}l@{\quad}l}
+  & сложение  & //  & целочисленное деление\\
-  & вычитание & \%  & остаток от деления\\
*  & умножение & **  & возведение в степень\\
/  & деление   &     &
\end{tabular}
\item \textbf{скобок}: \ $\dfrac{a+b}{c+d}\ \Rightarrow\ $ \texttt{(a + b) / (c + d)};
\item \textbf{переменных};
\item \textbf{числовых констант};
\item \textbf{математических функций}:\\[2pt]
\begin{tabular}{@{\qquad}l@{\qquad\qquad}l}
1 способ: & 2 способ:\\
\texttt{import math} & \texttt{from math import *}\\
\texttt{math.log(4.8)} & \texttt{log(4.8)}
\end{tabular}
\item $\sqrt x$~--- \texttt{sqrt(x)}; \ \texttt{x ** (1/2)}~--- \textbf{комплексное число!} (если ввели
отрицательное число);
\item \textbf{тригонометрических функций}~--- аргумент в радианах;
\item \texttt{degrees}~--- перевод из радиан в градусы; \ \texttt{radians}~--- перевод $x$ из градусов
в радианы.
\end{enumerate}

\subsection{Приоритет операций}
\begin{enumerate}[label=\arabic*)]
\item то, что в скобках \texttt{( )};
\item математические функции;
\item[3--4)] операция возведения в степень; унарная операция \texttt{(-)}: \texttt{-5},
\texttt{-sin(x)}, \texttt{+5} (в мат.~--- <<комплементарная>>);
\item[5)] \texttt{*} \ \texttt{/} \ \texttt{\%} \ \texttt{//};
\item[6)] \texttt{+} \ \texttt{-}.
\end{enumerate}

\texttt{x = 1 + 2 + 3}~--- по очереди слева направо.

\texttt{2 ** 3 ** 2} $\to$ \texttt{512}~--- справа налево: $2^{(3^2)}=2^9$.

\texttt{\#}~--- однострочный комментарий; \ \texttt{'''} многострочный комментарий \texttt{'''}.

% =====================================================================
\section{Линейные алгоритмы}
% =====================================================================
\begin{opr}
\textbf{Определение.} \emph{Алгоритм}~--- это конечная последовательность действий решения
какой-либо задачи.
\end{opr}

\lead{Способы представления алгоритма:}
\begin{enumerate}[label=\arabic*.]
\item структурограммы;
\item схема алгоритма;
\item псевдокод;
\item словесная запись.
\end{enumerate}

\subsection{Основные блоки схемы алгоритма}
Размеры даны в клетках: обычный блок~--- 6 клеток в ширину и 2 в высоту.
\begin{center}
\begin{tikzpicture}[sa]
\tikzset{lbl/.style={font=\small,align=center,text width=4.8cm,anchor=north},
         dim/.style={font=\scriptsize,text=gray}}
% ряд 1
\node[proc] (b1) at (0,0) {};
\node[term] (b2) at (5.6,0) {};
\node[io]   (b3) at (11.2,0) {};
\node[dim,above] at (b1.north) {6};  \node[dim,left] at (b1.west) {2};
\node[lbl] at (0,-0.75)    {1) вычислительный процесс};
\node[lbl] at (5.6,-0.75)  {2) начало и конец процесса};
\node[lbl] at (11.2,-0.75) {3) ввод и вывод};
% ряд 2
\node[dec]  (b4) at (0,-2.9) {};
\node[prep] (b5) at (5.6,-2.9) {};
\node[lpbeg,minimum width=2.4cm,minimum height=0.8cm] (b6a) at (10.1,-2.9) {};
\node[lpend,minimum width=2.4cm,minimum height=0.8cm] (b6b) at (12.7,-2.9) {};
\node[dim,above] at (b4.north) {6};  \node[dim,left] at (b4.west) {2};
\node[dim,above] at (b5.north) {6};  \node[dim,left] at (b5.west) {2};
\node[lbl] at (0,-3.75)    {4) решение};
\node[lbl] at (5.6,-3.75)  {5) цикл с параметром\\(старый ГОСТ)};
\node[lbl] at (11.4,-3.75) {6) цикл с параметром (новый ГОСТ):\\начало и конец цикла};
% ряд 3
\draw (0,-6.2) circle (0.5);
\node[dim] at (0,-5.5) {2};  \node[dim] at (-0.7,-6.2) {2};
\draw (5.2,-5.8)--(6.0,-5.8)--(6.0,-6.35)--(5.6,-6.65)--(5.2,-6.35)--cycle;
\node[dim] at (5.6,-5.6) {1};  \node[dim] at (5.0,-6.2) {1};
\node[pred] (b9) at (11.2,-6.2) {};
\node[dim,above] at (b9.north) {6};  \node[dim,left] at (b9.west) {2};
\node[lbl] at (0,-6.95)    {7) соединитель};
\node[lbl] at (5.6,-6.95)  {8) межстраничный\\соединитель};
\node[lbl] at (11.2,-6.95) {9) предопределённый процесс};
\end{tikzpicture}
\end{center}

\subsection{Правила составления схем}
\begin{enumerate}[label=\arabic*)]
\item Одинаковая ширина блоков~--- 6--8 клеток.
\item Высота блоков может быть различной.
\item Соединительная стрелка между блоками~--- 1 клетка, по необходимости~2.
\item Математические функции пишем как в математике.
\item Подряд идущие ввод данных, вывод данных, присваивания объединяем в один блок.
\item Каждое присваивание в блоке пишем с новой строки.
\item Схема алгоритма изображается без привязки к конкретному языку программирования.
\item Схема алгоритма должна занимать в высоту не более листа А4.
\end{enumerate}

\subsection{Подпрограммы}
В программировании часто встречаются алгоритмы, в которых повторяются фрагменты, одинаковые
по выполняемым действиям и различающиеся только значениями обрабатываемых данных.
\begin{opr}
\textbf{Определение.} Такие фрагменты записывают отдельно, однократно и называют
\emph{подпрограммами} (п/п). В дальнейшем к п/п можно обращаться многократно с различными
значениями данных.
\end{opr}

\lead{Синтаксис описания п/п}
(\texttt{[ ]}~--- часть, которая может отсутствовать):
\begin{code}
def <имя п/п>([<список формальных параметров>]):
    <команды>
    [return <возвращаемые значения>]
\end{code}

\lead{Синтаксис вызова п/п:}
\begin{code}
<имя п/п>([<список фактических параметров>])
\end{code}

\begin{opr}
\textbf{Определение.} \emph{Формальные параметры}~--- это имена переменных, которые используются
для описания п/п.

\emph{Фактические параметры}~--- это конкретные значения параметров, с которыми п/п будет вызвана
и выполнена.
\end{opr}

П/п бывают двух видов: \cn{1} функции; \ \cn{2} процедуры.
\begin{opr}
\textbf{Определение.} \emph{Функцией} будем считать п/п, в которой присутствует \texttt{return}
и возвращаемое значение. П/п без \texttt{return}~--- \emph{процедура}.
\end{opr}

\prim Процедура без параметров:
\begin{code}
def Error():
    print('Ошибка!')

a = 5
b = 0
if b == 0:
    Error()
else:
    print(a / b)
\end{code}

\prim Процедура с параметрами:
\begin{code}
def hello(name):
    print('Привет,', name, '!')

name = input('Вв. своё имя: ')
hello(name)
hello('Мама')
hello(1)
\end{code}

\prim Функция сложения:
\begin{code}
def func(a, b):
    rez = a + b
    return rez

z = func(6, 4)
print(z)            # 10
\end{code}

Процедура и функция различаются при вызове. В место вызова функции возвращается значение,
которое можно присваивать, выводить, использовать в других выражениях и~т.\,д.

\lead{Правила соответствия между формальными и фактическими параметрами:}
\begin{enumerate}[label=\cn{\arabic*}]
\item В классическом случае число формальных и фактических параметров должно совпадать.
\item Типы параметров (всех: и формальных, и фактических) в общем случае должны совпадать.
\item Порядок параметров (всех) нужно соблюдать.
\item Имена формальных и фактических параметров могут совпадать.
\end{enumerate}

\begin{opr}
\textbf{Определение.} \emph{Глобальная переменная}~--- это переменная, область видимости которой~---
вся программа. В Питоне это переменные, созданные в основной программе.

\emph{Локальная переменная}~--- это переменная, область видимости которой чем-то ограничена
(например, телом п/п).
\end{opr}

% =====================================================================
\section{Применение функций пользователя}
% =====================================================================
\begin{opr}
\textbf{Определение.} \emph{Необязательные параметры.} При описании п/п у формальных параметров
можно указать их значения по умолчанию. Тогда эти параметры при вызове можно не указывать, т.\,е.
они являются необязательными. Параметры, у которых значения по умолчанию нет, будут обязательными.
\end{opr}
\begin{code}
def proc(arg1, arg2=5, arg3=88):
    print(arg1, arg2, arg3)

proc(1)               # 1 5 88
proc(1, 2)            # 1 2 88
proc(888, arg3=77)    # 888 5 77
\end{code}

\primn{3.1} Программа вычисления суммы площадей двух прямоугольников\ispr{В рукописи стороны
второго прямоугольника тоже вводятся в \texttt{a1}, \texttt{b1}; должно быть \texttt{a2}, \texttt{b2}~---
иначе \texttt{sq(a2, b2)} даст ошибку.}.

Сразу после заголовка функции в тройных кавычках пишут её описание:
1)~что делает; 2)~перечисляют все входные значения; 3)~описывают, что возвращает \texttt{return}.
\begin{code}
# Программа вычисления суммы площадей двух прямоугольников
# Переменные:
#     a1, b1, a2, b2 - стороны прямоугольников
#     S1, S2 - их площади

# import ...  (сначала подключаются модули)

def sq(a, b):
    ''' Функция вычисления площади прямоугольника S
    Входные параметры:
        a, b - стороны прямоугольника
    Выходные данные: S
    '''
    S = a * b
    return S            # или сразу: return a * b

print('Вв. стороны 1-го прям.')
a1 = float(input())
b1 = float(input())
print('Вв. стороны 2-го прям.')
a2 = float(input())
b2 = float(input())
# эхо-печать
print(a1, b1, a2, b2)

S1 = sq(a1, b1)
S2 = sq(a2, b2)
print('Пл. 1-го прям. =', S1)
print('Пл. 2-го прям. =', S2)
print('Их сумма', S1 + S2)
\end{code}

\subsection{Правила оформления схем алгоритмов программ с п/п}
\begin{enumerate}[label=\arabic*)]
\item Начало и конец пишут только в основной программе. В п/п пишут \emph{вход} и \emph{выход}.
\item Для каждой п/п изображается отдельная схема алгоритма.
\item В верхнем блоке п/п нужно указать входные параметры, а в нижнем~--- выходные параметры
или вычисляемые значения.
\item Вызов п/п в основной программе изображается блоком <<предопределённый процесс>>
\tikz[sa,baseline=-0.6ex]\node[pred,minimum width=1.4cm,minimum height=0.45cm]{};.
\end{enumerate}

\begin{schema}
Схемы к примеру~3.1: слева~--- п/п \texttt{Sq}, справа~--- основная программа.
\begin{center}
\begin{tikzpicture}[sa]
% п/п Sq
\node[term] (s1) at (0,0) {Вход в п/п\\$Sq(a,b)$};
\node[proc,below=6mm of s1] (s2) {$S=a\cdot b$};
\node[term,below=6mm of s2] (s3) {Выход из п/п\\$Sq(S)$};
\draw (s1)--(s2)--(s3);
\gostcomment{s1.east}{$a,b$}
\gostcomment{s3.east}{$S$}
% основная программа
\node[term] (m1) at (8,0.6) {начало};
\node[io,below=of m1] (m2) {Ввод\\$a_1,b_1,a_2,b_2$};
\node[pred,below=of m2] (m3) {$S_1=Sq(a_1,b_1)$};
\node[pred,below=of m3] (m4) {$S_2=Sq(a_2,b_2)$};
\node[io,below=of m4] (m5) {Вывод\\$S_1,\ S_2,\ S_1+S_2$};
\node[term,below=of m5] (m6) {конец};
\draw (m1)--(m2)--(m3)--(m4)--(m5)--(m6);
\end{tikzpicture}
\end{center}
\end{schema}

% =====================================================================
\section{Условный оператор}
% =====================================================================
Условный оператор позволяет выполнить один из входящих в него наборов команд.

\lead{Синтаксис:}
\begin{code}
if <выражение>:
    <команда 1>
[else:
    <команда 2>]
\end{code}
Сначала проверяется значение выражения после \texttt{if}. Если оно истинно, то выполнится
\texttt{<команда 1>}, иначе (значение выражения ложно) выполнится \texttt{<команда 2>} (если есть
\texttt{else}).

\begin{minipage}[t]{0.48\textwidth}
\lead{За истину в Питоне считается:}
\begin{itemize}
\item \texttt{True};
\item числовое значение, отличное от нуля;
\item непустые объекты.
\end{itemize}
\end{minipage}\hfill
\begin{minipage}[t]{0.48\textwidth}
\lead{Ложь:}
\begin{itemize}
\item \texttt{False};
\item \texttt{0} (целое, вещественное \texttt{0.0});
\item пустые объекты (\texttt{None}).
\end{itemize}
\end{minipage}

\subsection{Схема условного оператора}
\begin{center}
\begin{tikzpicture}[sa,bx/.style={proc,minimum width=2.6cm}]
% полная форма
\node[dec] (d) at (0,0) {условие};
\draw (0,1.1) -- (d.north);
\node[bx] (c2) at (-2.2,-1.6) {команда 2};
\node[bx] (c1) at (2.2,-1.6) {команда 1};
\draw[->] (d.west) -| node[yes,above,pos=0.25]{нет} (c2.north);
\draw[->] (d.east) -| node[yes,above,pos=0.25]{да} (c1.north);
\draw (c2.south) |- (0,-2.7) -| (c1.south);
\draw[->] (0,-2.7) -- (0,-3.3);
\node[font=\small] at (0,-4.0) {полная форма};
% сокращённая форма, вариант 1
\begin{scope}[xshift=7.2cm]
\node[dec] (d) at (0,0) {условие};
\draw (0,1.1) -- (d.north);
\node[bx] (c1) at (2.2,-1.6) {команда 1};
\draw[->] (d.east) -| node[yes,above,pos=0.25]{да} (c1.north);
\draw (d.west) -- node[yes,above]{нет} (-2.0,0) |- (0,-2.7) -| (c1.south);
\draw[->] (0,-2.7) -- (0,-3.3);
\node[font=\small,align=center] at (0.6,-4.0) {сокращённая форма};
\end{scope}
\end{tikzpicture}
\end{center}

\begin{schema}
\lead{Условный оператор с одним условием (сокращённая форма)} можно изобразить и так:
\begin{center}
\begin{tikzpicture}[sa]
\node[dec] (d) at (0,0) {условие};
\draw (0,1.1) -- (d.north);
\node[proc,below=7mm of d] (c1) {команда 1};
\draw[->] (d.south) -- node[yes,right]{да} (c1.north);
\draw (d.west) -- node[yes,above]{нет} (-2.1,0) |- (0,-2.75);
\draw (c1.south) -- (0,-2.75);
\draw[->] (0,-2.75) -- (0,-3.35);
\end{tikzpicture}
\end{center}
\end{schema}

\subsection{Логические выражения}
В классическом варианте выражение после \texttt{if} должно быть \emph{логическим} (истина или ложь).
Чтобы выражение было логическим, оно должно включать:
\begin{itemize}
\item \textbf{операции отношения (сравнения)}: \ \texttt{>}, \ \texttt{<}, \ \texttt{>=}, \ \texttt{<=},
\ \texttt{==}, \ \texttt{!=};
\item \textbf{логические операции}:\\
\texttt{and}~--- логическое <<и>> (конъюнкция);\\
\texttt{or}~--- логическое <<или>> (дизъюнкция);\\
\texttt{not}~--- логическое <<не>> (отрицание).
\end{itemize}

\textbf{<<Исключающее или>>}~--- побитовое сложение по модулю~2, обозначается \texttt{\^{}}. Истина,
когда аргументы имеют противоположные значения; ложь~--- когда одинаковые.

\imp В Питоне у операций сравнения приоритет выше, чем у логических операций.

\prim $1\le x\le3$.\\
Классический вариант: \ \texttt{x >= 1 and x <= 3}.\\
В Питоне реализованы <<ленивые>> цепочки операций сравнения: \ \texttt{1 <= x <= 3}.
Сначала вычисляется \texttt{1 <= x}. Если это \texttt{False}, то дальше не идём (результат
\texttt{False}). Если \texttt{True}, то берём~\texttt{x} и дальше сравниваем: \texttt{x <= 3}.

\subsection{Тернарный условный оператор}
Конструкцию
\begin{code}
if X:
    A = Y
else:
    A = Z
\end{code}
можно записать в одну строку\dop{В рукописи под заголовком записана только полная форма
\texttt{if}--\texttt{else}; строка с тернарным оператором добавлена.}:
\begin{code}
A = Y if X else Z
\end{code}

\prim Найти квадрат максимума из трёх значений.

Переменные: \texttt{a, b, c}~--- исходные числа; \ \texttt{z}~--- максимум из них.
\begin{code}
print('введите 3 числа')
a = float(input())
b = float(input())
c = float(input())
# эхо-печать
print(a, b, c)
if a > b:
    z = a
else:
    z = b
if c > z:
    z = c
print("Квадрат макс =", z * z)
\end{code}

\begin{center}
\begin{tikzpicture}[sa,bx/.style={proc,minimum width=2.6cm}]
\node[term] (b) at (0,0) {начало};
\node[io] (in) at (0,-1.5) {Ввод $a,b,c$};
\node[io] (o1) at (0,-3.0) {Вывод $a,b,c$};
\node[dec] (d1) at (0,-4.6) {$a>b$};
\node[bx] (zb) at (-2.4,-6.0) {$z=b$};
\node[bx] (za) at (2.4,-6.0) {$z=a$};
\node[dec] (d2) at (0,-8.0) {$c>z$};
\node[bx] (zc) at (2.4,-9.4) {$z=c$};
\node[io] (o2) at (0,-11.6) {Вывод $z^2$};
\node[term] (e) at (0,-13.1) {конец};
\draw (b)--(in)--(o1)--(d1);
\draw[->] (d1.west) -| node[yes,above,pos=0.25]{нет} (zb.north);
\draw[->] (d1.east) -| node[yes,above,pos=0.25]{да} (za.north);
\draw (zb.south) |- (0,-6.9) -| (za.south);
\draw[->] (0,-6.9) -- (d2.north);
\draw[->] (d2.east) -| node[yes,above,pos=0.25]{да} (zc.north);
\draw (d2.west) -- node[yes,above]{нет} (-2.4,-8.0) |- (0,-10.4) -| (zc.south);
\draw[->] (0,-10.4) -- (o2.north);
\draw (o2)--(e);
\end{tikzpicture}
\end{center}

\primn{4.2} (к лаб.~4.1). Вычисление выражения с условием:
\[
y=\begin{cases}
\frac12\sqrt x, & x\ge1,\\[2pt]
\frac13\sqrt[3]{x}, & 0<x<1,\\[2pt]
\frac14\sqrt[4]{|x|}, & -1\le x\le0,\\[2pt]
0 & \text{иначе}.
\end{cases}
\]
\begin{code}
# Вычисление выражения с условием
# Переменные:
#     x - аргумент
#     y - вычисляемая функция
from math import sqrt

x = float(input('введите x \n'))

if x >= 1:
    y = 0.5 * sqrt(x)
else:
    if x > 0 and x < 1:
        y = 1/3 * x ** (1/3)
    else:
        if x >= -1 and x <= 0:
            y = 1/4 * abs(x) ** (1/4)
        else:
            y = 0

print('x =', x, 'y =', y)
\end{code}

\begin{center}
\begin{tikzpicture}[sa,bx/.style={proc,minimum width=2.6cm},dd/.style={dec,minimum width=2.6cm}]
\node[term] (b) at (1.5,0) {начало};
\node[io] (in) at (1.5,-1.5) {Ввод $x$};
\node[dd] (d1) at (1.5,-3.1) {$x\ge1$};
\node[bx] (y1) at (5.3,-4.6) {$y=\frac12\sqrt x$};
\node[dd] (d2) at (-1.3,-4.6) {$0<x<1$};
\node[bx] (y2) at (2.0,-6.1) {$y=\frac13\sqrt[3]{x}$};
\node[dd] (d3) at (-4.1,-6.1) {$x\ge-1$};
\node[bx] (y3) at (-1.1,-7.6) {$y=\frac14\sqrt[4]{|x|}$};
\node[bx] (y4) at (-6.4,-7.6) {$y=0$};
\node[io] (o) at (-0.5,-9.9) {Вывод $x,\ y$};
\node[term] (e) at (-0.5,-11.4) {конец};
\draw (b)--(in)--(d1);
\draw[->] (d1.east) -| node[yes,above,pos=0.2]{да} (y1.north);
\draw[->] (d1.west) -| node[yes,above,pos=0.2]{нет} (d2.north);
\draw[->] (d2.east) -| node[yes,above,pos=0.25]{да} (y2.north);
\draw[->] (d2.west) -| node[yes,above,pos=0.25]{нет} (d3.north);
\draw[->] (d3.east) -| node[yes,above,pos=0.25]{да} (y3.north);
\draw[->] (d3.west) -| node[yes,above,pos=0.25]{нет} (y4.north);
\draw (y4.south) |- (5.3,-8.8) (y3.south) -- (-1.1,-8.8) (y2.south) -- (2.0,-8.8) (y1.south) -- (5.3,-8.8);
\draw[->] (-0.5,-8.8) -- (o.north);
\draw (o)--(e);
\end{tikzpicture}
\end{center}

\zam Если функция где-то не определена, нужно разделить числовую прямую на части и для каждой
части записать свою ветку.

\primn{4.3} \ $y=\begin{cases}2x^2, & 0\le x\le2,\\ 3x^3, & x>2;\end{cases}$ \ при $x<0$ функция не определена\dop{Условия в рукописи нет, оно восстановлено по программе.}.
\begin{code}
# пример 4_3
x = float(input())
if x < 0:
    print('x =', x, 'ф. не определена')
else:
    if x <= 2:
        y = 2 * x ** 2
    else:
        y = 3 * x ** 3
    print('x =', x, 'y =', y)
\end{code}

\begin{center}
\begin{tikzpicture}[sa]
\node[term] (b) at (0,0) {начало};
\node[io,below=of b] (in) {Ввод $x$};
\node[dec,below=6mm of in] (d1) {$x<0$};
\node[io,minimum width=3.4cm] (er) at (4.6,-4.3) {Вывод $x$,\\<<ф. не опред.>>};
\node[dec] (d2) at (-3,-4.3) {$x\le2$};
\node[proc,minimum width=2.6cm] (y3) at (-5.4,-6.2) {$y=3x^3$};
\node[proc,minimum width=2.6cm] (y2) at (-0.6,-6.2) {$y=2x^2$};
\node[io] (out) at (-3,-8.3) {Вывод $x,\ y$};
\node[term] (e) at (0,-10.5) {конец};
\draw (b)--(in)--(d1);
\draw[->] (d1.east) -| node[yes,above,pos=0.2]{да} (er.north);
\draw[->] (d1.west) -| node[yes,above,pos=0.2]{нет} (d2.north);
\draw[->] (d2.west) -| node[yes,above,pos=0.3]{нет} (y3.north);
\draw[->] (d2.east) -| node[yes,above,pos=0.3]{да} (y2.north);
\draw (y3.south) |- (-3,-7.25) -| (y2.south);
\draw[->] (-3,-7.25) -- (out.north);
\draw (out.south) |- (0,-9.5);
\draw (er.south) |- (0,-9.5);
\draw[->] (0,-9.5) -- (e.north);
\end{tikzpicture}
\end{center}

\primn{4.4} (к лаб.~4.3). Определить, попадает ли точка $(x,y)$ в заштрихованную область:
\begin{center}
\begin{tikzpicture}[scale=0.75,line width=0.6pt]
\fill[pattern=north east lines] (0,-2) arc (-90:-270:2) -- (2,2) -- (2,-2) -- cycle;
\draw (0,-2) arc (-90:-270:2) -- (2,2) -- (2,-2) -- cycle;
\draw[->] (-2.8,0)--(3,0) node[right] {$x$};
\draw[->] (0,-2.8)--(0,2.8) node[above] {$y$};
\node[above left] at (0,2) {$2$};   \node[below left] at (0,-2) {$-2$};
\node[above left] at (-2,0) {$-2$}; \node[below right] at (2,0) {$2$};
\end{tikzpicture}
\end{center}
\lead{Подсказка:} отдельно записать условие для $x<0$, где $x^2+y^2\le R^2$.

% =====================================================================
\section{Многоальтернативное ветвление}
% =====================================================================
В Питоне \texttt{if}, вложенный в ветку \texttt{else}, функционально заменяется на \texttt{elif}.

\lead{Синтаксис}\ispr{В рукописи после первой строки пропущены \texttt{<команды 1>}.}
(\texttt{[ ]}~--- часть, которая может отсутствовать):
\begin{code}
if <выражение 1>:
    <команды 1>
[elif <выражение 2>:
    <команды 2>]
...
[else:
    <команды N+1>]
\end{code}

\primn{5.1} (к лаб.~5). Соотношение между мерами длины.

Переменные: \texttt{k}~--- код пункта меню.
\begin{code}
print('1 - миля')
print('2 - верста\n3 - сажень')
print('''4 - аршин
5 - фут''')
k = int(input('Вв. код пункта меню: '))
if k == 1:
    print('1,61 км')
elif k == 2:
    print('1,07 км')
elif k == 3:
    print('2,13 м')
elif k == 4:
    print('0,71 м')
elif k == 5:
    print('0,30 м')
else:
    print('Нет такого пункта меню')
\end{code}
\vspace{-4pt}
{\small В рукописи после ветки \texttt{k == 2} стоит многоточие; ветки для \texttt{k = 3, 4, 5} и
\texttt{else} дописаны\dop{1 сажень $=2{,}1336$~м, 1 аршин $=0{,}7112$~м, 1 фут $=0{,}3048$~м.}.}

\subsection{Оператор выбора \texttt{match}}
\lead{Синтаксис:}
\begin{code}
match <выражение-селектор>:
    case <шаблон 1>:
        <команды 1>
    ...
    case <шаблон N>:
        <команды N>
    [case _:
        <команды N+1>]
\end{code}
\texttt{case \_}~--- ветка <<иначе>>: выполняется, если не подошёл ни один шаблон. Несколько
шаблонов в одной ветке перечисляют через \texttt{|}.

\primn{5.3} Простейший калькулятор.
\begin{code}
# Переменные:
#     a, b - операнды
#     y - результат
#     znak - знак операции
#     flag - признак существования знака операции
print('?')
a = float(input())
znak = input()
b = float(input())
flag = True
match znak:
    case '+':
        y = a + b
    case '-':
        y = a - b
    case '*':
        y = a * b
    case '/' | ':':
        y = a / b
    case _:
        flag = False
if flag:
    print(a, znak, b, '=', y)
else:
    print('Неизвестный знак операции')
\end{code}
\vspace{-4pt}
{\small Последние четыре строки (вывод результата) дописаны: в рукописи программа и схема
обрываются после \texttt{match}\dop{Добавлены вывод результата и блок <<flag>> на схеме.}.}

\begin{center}
\begin{tikzpicture}[sa,
  bx/.style={proc,minimum width=2.6cm}]
\node[term] (b) at (0,0) {начало};
\node[io] (in) at (0,-1.5) {Ввод $a$, знак, $b$};
\node[proc] (fl) at (0,-3.0) {flag $=$ True};
\node[dec] (z) at (0,-4.7) {знак};
\draw (b)--(in)--(fl)--(z);
\draw (z.south) -- (0,-5.8);
\draw (-6.8,-5.8) -- (6.8,-5.8);
\foreach \xx/\txt/\lab/\nm in {-6.8/{$y=a+b$}/{'+'}/c1, -3.4/{$y=a-b$}/{'$-$'}/c2,
    0/{$y=a\cdot b$}/{'*'}/c3, 3.4/{$y=a/b$}/{'/', ':'}/c4, 6.8/{flag $=$ False}/{иначе}/c5}
  { \node[bx] (\nm) at (\xx,-6.9) {\txt};
    \draw[->] (\xx,-5.8) -- node[yes,left,pos=0.4]{\lab} (\nm.north); }
\draw (-6.8,-7.4) -- (-6.8,-7.9) -- (6.8,-7.9) -- (6.8,-7.4);
\foreach \xx in {-3.4,0,3.4} \draw (\xx,-7.4) -- (\xx,-7.9);
\node[dec] (f) at (0,-9.0) {flag};
\draw[->] (0,-7.9) -- (f.north);
\node[io,minimum width=3.4cm] (o1) at (3.8,-10.7) {Вывод\\$a$, знак, $b$, $y$};
\node[io,minimum width=3.4cm] (o2) at (-3.8,-10.7) {Вывод <<Неизвестный\\знак операции>>};
\draw[->] (f.east) -| node[yes,above,pos=0.2]{да} (o1.north);
\draw[->] (f.west) -| node[yes,above,pos=0.2]{нет} (o2.north);
\node[term] (e) at (0,-12.8) {конец};
\draw (o1.south) |- (0,-11.9) (o2.south) |- (0,-11.9);
\draw[->] (0,-11.9) -- (e.north);
\end{tikzpicture}
\end{center}

\begin{schema}
\lead{Общий вид выбора на схеме.} От блока <<решение>> идёт одна линия, от неё ответвляются
выходы, подписанные значениями селектора; после блоков ветви сходятся:
\begin{center}
\begin{tikzpicture}[sa]
\node[dec,minimum width=2.2cm] (k) at (0,0) {$k$};
\draw (k.south) -- (0,-5.4);
\node[proc,minimum width=2.6cm] (p1) at (2.6,-1.5) {};
\node at (2.6,-2.7) {$\dots$};
\node[proc,minimum width=2.6cm] (p3) at (2.6,-3.9) {};
\draw[->] (0,-1.5) -- node[yes,above]{1} (p1.west);
\draw[->] (0,-2.7) -- node[yes,above]{2} (1.1,-2.7);
\draw[->] (0,-3.9) -- node[yes,above]{$N$} (p3.west);
\draw (p1.east) -- (4.4,-1.5) (p3.east) -- (4.4,-3.9);
\draw[->] (4.4,-1.5) -- (4.4,-4.7) -- (0,-4.7);
\end{tikzpicture}
\end{center}
\end{schema}

% =====================================================================
\section{Циклы с параметром (со счётчиком) \texttt{for}}
% =====================================================================
Когда одни и те же действия необходимо выполнить многократно, для этого используют различные
операторы циклов. В программировании различают следующие \textbf{типы циклов}:
\begin{enumerate}[label=\arabic*)]
\item циклы с предусловием;
\item циклы с постусловием;
\item циклы с параметром;
\item циклы с перебором коллекций.
\end{enumerate}

Каждый из перечисленных циклов имеет \textbf{общие структурные особенности} и состоит из 4 частей:
\begin{enumerate}[label=\arabic*)]
\item \emph{подготовка цикла}~--- необходима для задания начальных значений;
\item \emph{тело цикла}~--- те команды, которые расположены внутри цикла и выполняются многократно;
\item \emph{модификация}~--- изменение параметра цикла;
\item \emph{проверка условия} окончания или продолжения цикла.
\end{enumerate}

\lead{Синтаксис цикла с параметром:}
\begin{code}
for <параметр> in <список значений>:
    <команды>
\end{code}
\lead{Как работает:} параметр поочерёдно принимает значения из списка, при этом каждый раз
выполняются команды из тела цикла.

\subsection{Диапазон \texttt{range()}}
Функция \texttt{range()} возвращает диапазон целых чисел от заданного начального значения до
конечного с заданным шагом:
\begin{center}
\texttt{range(<нач. знач.>, <след. за конечным знач.>, <шаг>)}
\end{center}
Если начальное значение не указано, то оно $=0$; если не указан шаг, то он $=+1$.

\begin{minipage}[t]{0.48\textwidth}
\begin{code}
for i in range(5):
    print(i)
\end{code}
\end{minipage}\hfill
\begin{minipage}[t]{0.48\textwidth}
\begin{code}
lst = [1, 2, 3]
for el in lst:
    print(el)
\end{code}
\end{minipage}

\begin{schema}
\lead{Схема цикла по старому ГОСТу}~--- блок <<модификация>> (шестиугольник), линия возврата
приходит в его левую вершину, выход из цикла~--- из правой:
\begin{center}
\begin{tikzpicture}[sa]
\node[prep] (h) at (0,0) {$i=0(1)4$};
\node[io,below=of h] (o) {Вывод $i$};
\cycleold{h}{o}
\draw[->] (h-out.center) -- ++(0,-0.5);
\node[prep] (h2) at (7,0) {$el$ из $lst$};
\node[io,below=of h2] (o2) {Вывод $el$};
\cycleold{h2}{o2}
\draw[->] (h2-out.center) -- ++(0,-0.5);
\end{tikzpicture}
\end{center}
\end{schema}

\begin{schema}
\lead{По новому ГОСТу} цикл изображается двумя блоками: начало цикла (с законом изменения
параметра) и конец цикла (с именем параметра):
\begin{center}
\begin{tikzpicture}[sa]
\node[lpbeg] (lb) {$i=0(1)4$};
\node[below=4mm of lb] (dots) {$\dots$};
\node[lpend,below=4mm of dots] (le) {$i$};
\draw (lb)--(dots)--(le);
\end{tikzpicture}
\end{center}
\end{schema}

\primn{6.1} Вывести $n_{\max}$ членов последовательности $a_1=2$, \ $a_n=n\cdot a_{n-1}$%
\dop{Условия в рукописи нет, оно восстановлено по программе.}.
\begin{code}
# Переменные:
#     nmax - количество чисел
#     n - параметр цикла
#     a1, a2 - предыдущее и текущее значения
nmax = int(input('Вв. кол-во чисел\n'))
a1 = 2
print('a1 = 2')
for n in range(2, nmax + 1):
    a2 = n * a1
    print('a%d = %d' % (n, a2))
    a1 = a2
\end{code}
\vspace{-4pt}
{\small Последняя строка \texttt{a1 = a2} добавлена\ispr{В рукописи нет \texttt{a1 = a2}: тогда
\texttt{a1} всегда равно~2 и вместо $a_n=n\,a_{n-1}$ печатается $2n$.}.}

\primn{6.2} Числа Фибоначчи: \ $f_n=f_{n-1}+f_{n-2}$, \ $f_1=1$, \ $f_2=1$.
\begin{code}
# Пример 6_2
# Числа Фибоначчи
# Переменные:
#     n - параметр цикла
#     nmax - количество чисел
#     f1, f2, f3 - предыдущие и текущее числа Фибоначчи
nmax = int(input('Вв. количество чисел: '))
f1 = 1
f2 = 1
print('f1 = 1\nf2 = 1')
for n in range(3, nmax + 1):
    f3 = f1 + f2
    print('f%d = %d' % (n, f3))
    f1 = f2
    f2 = f3
\end{code}

\begin{center}
\begin{tikzpicture}[sa]
\node[term] (b) {начало};
\node[io,below=of b] (in) {Ввод $n_{\max}$};
\node[proc,below=of in] (p) {$f_1=1$\\$f_2=1$};
\node[io,below=of p] (o) {Вывод $f_1,\ f_2$};
\node[prep,minimum width=3.4cm,below=6mm of o] (h) {$n=3(1)n_{\max}$};
\node[proc,below=of h] (s) {$f_3=f_1+f_2$};
\node[io,below=of s] (o2) {Вывод $n,\ f_3$};
\node[proc,below=of o2] (sh) {$f_1=f_2$\\$f_2=f_3$};
\node[term,below=14mm of sh] (e) {конец};
\draw (b)--(in)--(p)--(o)--(h)--(s)--(o2)--(sh);
\cycleold[0.6]{h}{sh}
\draw[->] (h-out.center) -- (e.north);
\end{tikzpicture}
\end{center}
{\small Схема из рукописи\ispr{В рукописи после блока $f_3=f_1+f_2$ стоит ещё блок ввода-вывода с той
же формулой; он лишний и убран.}.}

\subsection{Форматный вывод: оператор \texttt{\%}}
В строке ставят спецификаторы, а после строки через \texttt{\%} перечисляют значения:
\texttt{'f\%d = \%d' \% (n, f3)}.
\begin{code}
# %d - целое число
# %f - вещественное число в формате с фиксированной точкой
print('%f' % 1234568.12)     # 1234568.120000
print('%f' % 1e-3)           # 0.001000
# %e - вещественное число в формате с плавающей точкой
print('%e' % 1234568.12)     # 1.234568e+06
# %s - для вывода строк
print('%10d' % 123)          #        123   (поле из 10 позиций)
print('%8.2f' % 123.4567)    #   123.46     (8 позиций, 2 знака после точки)
print('%.3f' % 123.45678)    # 123.457
print('%20f' % 123.45678)    #           123.456780   (строка из 20 символов)
# %o - вывод числа в 8 СС
print('%o' % 123)            # 173
# %x - вывод целого в 16 СС
print('%x' % 123)            # 7b
\end{code}

\subsection{Накопление суммы}
\[
S=f_1+f_2+\dots+f_n=\sum_{i=1}^{n}f_i .
\]
Для вычисления такой суммы в программировании можно воспользоваться \emph{рекуррентной формулой}:
\[
S_i=S_{i-1}+f_i,\qquad i=1(1)n=\overline{1,n}\ \ (i=1,2,\dots,n),\qquad S_0=0.
\]
\imp Для того чтобы не испортить значение суммы, её начальное значение нужно задать равным нулю
($S_0=0$).

\primn{6.3} \ $S=\displaystyle\sum_{i=1}^{n}\frac1{i^2}=\frac11+\frac1{2^2}+\dots+\frac1{n^2}$.
\begin{code}
# Пример вычисления суммы
# Переменные:
#     i - параметр цикла
#     S - сумма
#     n - количество слагаемых
n = int(input("Вв. количество слагаемых: "))
S = 0
for i in range(1, n + 1):
    S = S + 1 / i ** 2      # или S += 1 / i ** 2
print('Сумма %d слагаемых = %f' % (n, S))
\end{code}

\subsection{Произведение}
\[
P=f_1\cdot f_2\cdot\ldots\cdot f_n=\prod_{i=1}^{n}f_i,\qquad
P_i=P_{i-1}\cdot f_i,\quad i=\overline{1,n},\qquad P_0=1.
\]

\primn{6.4} Натуральная степень действительного числа через произведение:
\[
a^n=\underbrace{a\cdot a\cdot\ldots\cdot a}_{n\ \text{раз}}=\prod_{i=1}^{n}a .
\]
\begin{code}
# Пример 6_4
# Вычисление через произведение
# натуральной степени заданного действительного числа
# Переменные:
#     P - произведение (результат)
#     a - основание
#     n - натуральная степень
#     i - параметр цикла
print('Введите основание и степень')
a = float(input())
n = int(input())
P = 1
for i in range(1, n + 1):    # можно range(n): i в теле не используется
    P *= a
print('%.2f в степени %d = %.3f' % (a, n, P))
\end{code}

\primn{6.5}
\[
z=\begin{cases}
\dfrac{x+1}{2}\displaystyle\sum_{n=1}^{10}\Bigl(\frac{x}{n}\Bigr)^{2}, & x\le2,\\[12pt]
\dfrac{\sin x+\cos x}{2+\sin x}\displaystyle\prod_{n=0}^{15}\Bigl(1+\frac{x}{n+2}\Bigr), & x>2.
\end{cases}
\]
\begin{code}
# Пример 6_5
# Переменные:
#     x - аргумент
#     z - вычисляемая функция
#     n - параметр цикла
from math import sin, cos
x = float(input('Введите x: '))
if x <= 2:
    S = 0
    for n in range(1, 10 + 1):
        S += (x / n) ** 2
    z = (x + 1) / 2 * S
else:
    P = 1
    for n in range(0, 16):
        P *= 1 + x / (n + 2)
    z = (sin(x) + cos(x)) / (2 + sin(x)) * P
print('x = %5.2f \t z = %5.2f' % (x, z))
\end{code}
\vspace{-4pt}
{\small Показатель степени в сумме~--- как в условии\ispr{В рукописи в условии $(x/n)^2$, а в программе
\texttt{(x/n)**n} и на схеме $(1/n)^n$. Здесь везде как в условии. Если на доске было $(x/n)^n$,
поменяй показатель во всех трёх местах.}.}

\begin{center}
\begin{tikzpicture}[sa,bx/.style={minimum width=3.4cm}]
\node[term] (b) at (0,0) {начало};
\node[io,below=of b] (in) {Ввод $x$};
\node[dec,below=6mm of in] (d) {$x\le2$};
% ветка "нет"
\node[proc,bx] (p0) at (-4,-4.4) {$P=1$};
\node[prep,bx,below=6mm of p0] (hp) {$n=0(1)15$};
\node[proc,bx,below=of hp] (pp) {$P=P\Bigl(1+\dfrac{x}{n+2}\Bigr)$};
\node[proc,bx,below=14mm of pp] (zp) {$z=\dfrac{\sin x+\cos x}{2+\sin x}\cdot P$};
\draw (p0)--(hp)--(pp);
\cycleold[0.5]{hp}{pp}
\draw[->] (hp-out.center) -- (zp.north);
% ветка "да"
\node[proc,bx] (s0) at (4,-4.4) {$S=0$};
\node[prep,bx,below=6mm of s0] (hs) {$n=1(1)10$};
\node[proc,bx,below=of hs] (ss) {$S=S+\Bigl(\dfrac{x}{n}\Bigr)^2$};
\node[proc,bx,below=14mm of ss] (zs) {$z=\dfrac{x+1}{2}\cdot S$};
\draw (s0)--(hs)--(ss);
\cycleold[0.5]{hs}{ss}
\draw[->] (hs-out.center) -- (zs.north);
\draw[->] (d.west) -| node[yes,above,pos=0.2]{нет} (p0.north);
\draw[->] (d.east) -| node[yes,above,pos=0.2]{да} (s0.north);
% слияние
\coordinate (m) at (0,-11.9);
\draw (zp.south) |- (m) (zs.south) |- (m);
\node[io,below=6mm of m] (o) {Вывод $x,\ z$};
\node[term,below=of o] (e) {конец};
\draw[->] (m)--(o); \draw (o)--(e);
\end{tikzpicture}
\end{center}

\subsection{Табулирование}
\begin{opr}
\textbf{Определение.} \emph{Протабулировать} функцию $y=f(x)$~--- это вычислить значения этой функции
при различных значениях аргумента.
\end{opr}
Закон изменения аргумента записывают так:
\[
x=x_0(h_x)x_n,\qquad x_0\text{~--- начальное значение},\ \ h_x\text{~--- шаг},\ \
x_n\text{~--- конечное значение аргумента}.
\]
\prim $x=1(0{,}5)3$, \ $y=f(x)=\sin x$:
\[
x=1:\ y=\sin1=\dots;\qquad x=1{,}5:\ y=\sin1{,}5=\dots;\qquad\dots;\qquad x=3:\ y=\sin3=\dots
\]

Поскольку в Питоне в цикле \texttt{for} при использовании функции \texttt{range()} можно работать
с параметрами цикла только целого типа, необходимо ввести переменную $N_x$, которая равна числу
повторений цикла при указанном законе изменения параметра~$x$:
\[
N_x=\Bigl[\frac{x_n-x_0}{h_x}\Bigr]+1,\qquad [\ \ ]\text{~--- целая часть.}
\]
\imp \texttt{for x in range(1, 3 + 0.5, 0.5)}~--- \textbf{нельзя!!!}

\primn{6.6} Протабулировать функцию $y=\sin x$.
\begin{code}
# Переменные:
#     x - аргумент
#     x0 - начальное значение аргумента x
#     xn - конечное значение аргумента x
#     hx - шаг изменения x
#     y - вычисляемая функция
#     Nx - количество точек в диапазоне
#     i - параметр цикла
from math import *
print('Введите x0(hx)xn')      # лучше подписать каждое значение
x0 = float(input())
hx = float(input())
xn = float(input())
Nx = trunc((xn - x0) / hx + 1e-6) + 1      # 1e-6 - погрешность
x = x0
for i in range(Nx):
    y = sin(x)
    print('x = %10.2f \t y = %10.3f' % (x, y))
    x = x + hx
\end{code}
\zam $10^{-6}$ добавляют на случай погрешности вещественной арифметики: например, при
$x=0(0{,}1)0{,}7$ в Питоне \texttt{0.7 / 0.1} даёт \texttt{6.999999999999999}, и без поправки
потерялась бы последняя точка.

% =====================================================================
\section{Цикл с предусловием \texttt{while}}
% =====================================================================
Цикл \texttt{for} используется тогда, когда заранее известно число повторений цикла. Когда оно
неизвестно, используют циклы с предусловием или с постусловием.

\lead{Синтаксис цикла с предусловием:}
\begin{code}
while <выражение>:
    <команды 1>
[else:
    <команды 2>]
\end{code}
Пока значение выражения истинно, выполняются \texttt{<команды 1>} из тела цикла. Как только
значение выражения становится ложным, происходит выход из цикла.

\imp При наличии \texttt{else} после \emph{естественного} выхода из цикла (для цикла с предусловием это
происходит тогда, когда выражение стало ложно) выполняются \texttt{<команды 2>}.

\begin{center}
\begin{tikzpicture}[sa]
\coordinate (top) at (0,0.9);
\node[dec] (d) at (0,0) {выражение};
\node[proc,below=7mm of d] (c) {команды};
\draw[->] (0,1.5) -- (d.north);
\draw[->] (d.south) -- node[yes,right]{да} (c.north);
\draw[->] (c.south) -- ++(0,-0.35) -| (-2.2,0.9) -- (0,0.9);
\draw (d.east) -- node[yes,above]{нет} (2.2,0) -- (2.2,-3.0) -- (0,-3.0);
\draw[->] (0,-3.0) -- (0,-3.6);
\end{tikzpicture}
\end{center}

\subsection{Цикл с постусловием}
Этот цикл отличается от цикла с предусловием тем, что проверка условия продолжения или окончания
цикла в нём происходит \emph{после} выполнения тела цикла.
\begin{center}
\begin{tikzpicture}[sa]
\node[proc] (c) at (0,0) {команды};
\node[dec,below=6mm of c] (d) {выражение};
\draw[->] (0,1.3) -- (c.north);
\draw (c)--(d);
\draw[->] (d.west) -- node[yes,above]{да/нет} ++(-0.9,0) |- (0,0.85);
\draw[->] (d.south) -- node[yes,right]{нет/да} ++(0,-0.7);
\end{tikzpicture}
\end{center}
Так как в цикле с предусловием условие проверяется до тела цикла, то, если условие сразу ложно,
цикл выполнится 0~раз. \textbf{Цикл с постусловием выполнится минимум 1~раз.}

\subsection{Общее для циклов \texttt{for} и \texttt{while}}
\lead{Оператор \texttt{continue}.} При выполнении \texttt{continue} происходит переход на следующий
шаг (итерацию) цикла; при этом на текущей итерации команды, стоящие в теле цикла после
\texttt{continue}, выполнены не будут.
\begin{code}
for i in range(1, 6):
    if i == 3:
        continue
    print('2^{} = {}'.format(i, 2 ** i))
\end{code}

\lead{Оператор \texttt{break}.} При выполнении \texttt{break} в теле цикла происходит досрочный выход
из цикла: на текущей итерации не выполняются команды тела цикла, стоящие после \texttt{break},
и не выполняются следующие итерации.
\begin{code}
for i in range(1, 6):
    if i == 3:
        break
    print(f'2^{i} = {2 ** i}')
\end{code}
\imp \texttt{break} и \texttt{continue} в лабах не используем (считаем, что их нет).

\lead{Конструкция \texttt{else}.} Ветка \texttt{else} выполняется, только если цикл закончился
естественно, а не через \texttt{break}:
\begin{code}
lst = [1, 2, 3, 4]
for el in lst:
    if el < 0:
        break
else:
    print('В списке ' + str(lst) + ' нет отрицательных')
\end{code}

\primn{7.1} (6.3 по \texttt{while}). Протабулировать функцию при $x=x_0(h_x)x_n$:
\[
y=\begin{cases}-x+2, & -1\le x\le2,\\ -3, & x<-1;\end{cases}
\qquad\text{при остальных $x$ функция не определена.}
\]
\begin{code}
# Табулирование циклом с предусловием
# Переменные:
#     x, x0, hx, xn - аргумент и параметры его изменения
#     y - функция
print('Введите x0(hx)xn')
x0 = float(input())
hx = float(input())
xn = float(input())
x = x0
while x <= xn + hx / 3:
    if x > 2:
        print(f'x = {x} функция не опр.')
    else:
        if x < -1:
            y = -3
        else:
            y = -x + 2
        print(f'x = {x} y = {y}')
    x = x + hx
\end{code}
\vspace{-4pt}
{\small Значение при $x<-1$~--- как в условии\ispr{В рукописи в условии $-3$, а в программе
\texttt{y = 3}. Если в условии опечатка, верни \texttt{y = 3}: при $y=3$ функция непрерывна в точке
$x=-1$.}. Добавка \texttt{hx / 3} в условии цикла~--- запас на погрешность, чтобы не потерять
точку $x_n$.}

\primn{7.2} \ $f_1=2+\sin3x$, \ $f_2=\cos\dfrac{x}{3}$. Определить аргумент, при котором расстояние
между функциями $f_1$ и $f_2$ на заданном интервале максимально.

Расстояние $d=\sqrt{(x_1-x_2)^2+(y_1-y_2)^2}$; \ при $x_1=x_2=x$ \
$d=\sqrt{(y_1-y_2)^2}=|y_1-y_2|$.
\begin{code}
# Пример 7_2
# Переменные:
#     x, x0, hx, xn - аргумент и его параметры
#     y1, y2 - функции
#     d - расстояние
#     dmax - максимальное расстояние
#     xmax - аргумент, при котором расстояние максимально
from math import sin, cos

def f1(x):
    '''Вычисляет и возвращает значение функции f1
    Формальный параметр: x - аргумент'''
    return 2 + sin(3 * x)

def f2(x):
    '''Вычисляет и возвращает значение функции f2'''
    return cos(x / 3)

print('Введите x0(hx)xn')
x0 = float(input())
hx = float(input())
xn = float(input())

x = x0
dmax = abs(f1(x0) - f2(x0))
xmax = x0
while x <= xn + hx / 4:
    y1 = f1(x)
    y2 = f2(x)
    d = abs(y1 - y2)
    print(f'{x = }\t f1 = {y1:10.2f} f2 = {y2:.3f} {d = :10.2f}')
    if d > dmax:
        dmax = d
        xmax = x
    x += hx
print(f'макс. расст. между ф. {dmax} в точке {xmax}')
\end{code}
\zam Начальное значение можно задать и так: \texttt{dmax = -inf} (\texttt{inf} есть в модуле
\texttt{math}) или \texttt{dmax = float('-inf')}\ispr{В рукописи \texttt{dmax = inf}, без минуса. При
$+\infty$ условие \texttt{d > dmax} не выполнится ни разу, и максимум не найдётся.}.

\zam Запись \texttt{\{x = \}} в f-строке печатает и имя, и значение: \texttt{x = 0.5}.

% =====================================================================
\section{Вложенные циклы}
% =====================================================================
\begin{opr}
\textbf{Определение.} \emph{Вложенным} называют цикл, содержащий внутри себя один или несколько
других циклов. Циклы, охватывающие другие циклы, называют \emph{внешними}, а остальные~---
\emph{внутренними}.
\end{opr}
\begin{enumerate}[label=\arabic*)]
\item Правила организации внешнего и внутреннего циклов такие же, как и для простого цикла
соответствующего вида.
\item Внутренний цикл должен полностью располагаться в теле внешнего цикла.
\end{enumerate}

Циклы, образующие вложенный цикл, условно разбивают на \emph{уровни} (\emph{ранги}): внешний цикл~---
уровень~0, первый внутренний цикл~--- уровень~1 и~т.\,д.

Параметры циклов разных уровней изменяются не одновременно. Сначала все свои значения
последовательно принимает параметр самого внутреннего цикла при неизменных значениях параметров
внешних для него циклов. После этого на один шаг изменяется параметр предыдущего по рангу цикла,
и снова все свои значения принимает параметр цикла последнего уровня вложения. Этот процесс
происходит до тех пор, пока параметры циклов всех уровней не примут все свои возможные значения.
\begin{code}
for i in range(1, 4):
    for j in range(3):
        print(i, j)
\end{code}
Если во вложенном цикле реализованы 2 циклические структуры с количеством повторений $N_1$ и $N_2$
соответственно, то общее число повторений тела внутреннего цикла равно $N_1\cdot N_2$ (в примере
выше $3\cdot3=9$ строк).

\primn{8.1} Табулирование функции двух переменных $z=f(x,y)$.

Переменные: \texttt{x, x0, xn, hx}~--- для изменения $x$; \ \texttt{y, y0, yn, hy}~--- для изменения
$y$; \ \texttt{z}~--- вычисляемая функция.

Пусть $z=\sin x\cdot\cos2y$, \ $x=x_0(h_x)x_n$, \ $y=y_0(h_y)y_n$.
\begin{code}
from math import sin, cos
print('Введите x0(hx)xn')
x0 = float(input())
hx = float(input())
xn = float(input())
print('Введите y0(hy)yn')
y0 = float(input())
hy = float(input())
yn = float(input())

x = x0
while x <= xn + hx / 3:
    y = y0
    while y <= yn + hy / 4:
        z = sin(x) * cos(2 * y)
        print(f'{x = } {y = } {z = }')
        y += hy
    x += hx
\end{code}
\imp Начальное значение параметра внутреннего цикла (\texttt{y = y0}) задаётся \emph{внутри} внешнего
цикла: перед каждым новым $x$ перебор $y$ начинается заново.

\primn{8.2}
\[
y=\frac{x+1{,}5}{3{,}75}\sum_{k=0}^{9}\ \prod_{n=5}^{3}(k+n)\cdot x^{k+n},\qquad x=x_0(h_x)x_n .
\]
\begin{code}
# Пример 8_2
# Переменные:
#     x, x0, xn, hx - параметры изменения x
#     y - вычисляемая функция
#     S, P - сумма и произведение
#     k, n - параметры циклов

# вводим x0, hx, xn   <- это обязательно

x = x0
while x <= xn + hx / 3:
    S = 0
    for k in range(10):
        P = 1
        for n in range(5, 2, -1):
            P *= (k + n) * x ** (k + n)
        S += P
    y = (x + 1.5) / 3.75 * S
    print(f'{x = :10.2f}\t{y = :10.3f}')
    x = x + hx
\end{code}
\zam Произведение по $n$ от~5 до~3 идёт в обратном порядке: \texttt{range(5, 2, -1)} даёт
$5,4,3$. Сумму обнуляют перед циклом по~$k$, произведение~--- перед циклом по~$n$, т.\,е. каждое
накопление начинают заново на своём уровне вложения.

\primn{8.3}
\[
y=\sum_{i=0}^{3}\ \prod_{j=5}^{2}\ \sum_{k=4}^{1}\frac{i+j\cdot k}{i\cdot j+k}.
\]
\begin{code}
# Переменные:
#     i, j, k - параметры циклов
#     P - произведение
#     S1, S2 - суммы
S1 = 0
for i in range(4):
    P = 1
    for j in range(5, 1, -1):
        S2 = 0
        for k in range(4, 0, -1):
            S2 += (i + j * k) / (i * j + k)
        P *= S2
    S1 += P
print('y =', S1)
\end{code}
\vspace{-4pt}
{\small Скобки в дроби расставлены по формуле\ispr{В рукописи \texttt{S2 += (i + j*k / (i*j + k))}:
по приоритету операций это $i+\frac{j\cdot k}{i\cdot j+k}$, а не дробь из условия.}.}

\primn{8.4}
\[
y=a\prod_{i=1}^{10}\ \prod_{j=1}^{i}\frac{i}{j+|a|^{j}},\qquad a=a_0(h_a)a_n .
\]
\begin{code}
# Переменные:
#     a, a0, ha, an - параметры переменной a
#     y - функция
#     i, j - параметры циклов
#     P1, P2 - произведения
#     P - произведение для вычисления |a|**j

# ввод a0, ha, an

a = a0
while a <= an + ha / 4:
    P1 = 1
    for i in range(1, 11):
        P2 = 1
        P = 1
        for j in range(1, i + 1):
            P = P * abs(a)
            P2 = P2 * i / (j + P)
        P1 = P1 * P2
    y = a * P1
    print(f'{a = } {y = }')
    a += ha
\end{code}
\zam Степень $|a|^j$ здесь не вычисляется через \texttt{**}, а накапливается произведением:
на каждом шаге внутреннего цикла \texttt{P} умножается на~$|a|$, поэтому при текущем~$j$
\texttt{P}~$=|a|^j$. Для этого \texttt{P = 1} задают перед внутренним циклом.

% =====================================================================
\section{Списки (одномерные массивы)}
% =====================================================================
\begin{opr}
\textbf{Определение.} \emph{Список} в Питоне~--- это индексированная (упорядоченная) динамическая
коллекция, которая состоит из одинаковых или различных элементов.

\emph{Динамическая}~--- это значит, что количество элементов может меняться во время выполнения
программы.
\end{opr}
\imp \emph{Массив}, в отличие от списка в Питоне, состоит из элементов одного типа.

\lead{Создание списка} (\texttt{list()}):
\begin{code}
# создание списка
lst1 = [4, 5, 1]
lst2 = [1, 'строка', 3.45]
# создание пустого списка
empt1 = []
empt2 = list()
\end{code}

Список~--- \emph{индексированная и упорядоченная} коллекция: у каждого элемента списка есть
порядковый номер (индекс), по которому можно обращаться к элементу. Нумерация с нуля;
отрицательный индекс считает с конца.
\begin{code}
a = [1, 2, 3, 4]
print(a[2])            # 3
print(a[-1])           # 4
print([4, 3, 6][2])    # 6
print(a[99])           # Ошибка! (IndexError)
\end{code}

\lead{Изменение элемента:}
\begin{code}
a[0] = 99
print(a)               # [99, 2, 3, 4]
\end{code}

\lead{Метод \texttt{append()}}~--- добавление в конец списка:
\begin{code}
a.append(13)
print(a)               # [99, 2, 3, 4, 13]
\end{code}

\lead{Метод \texttt{extend()}}~--- расширение одного списка элементами другого:
\begin{code}
a1 = [1, 2, 3]
a2 = [44, 55]
a1.extend(a2)
print(a1)              # [1, 2, 3, 44, 55]
\end{code}

\begin{code}
a = 1
b = 2
c = [a, b]
print(c)               # [1, 2]
b = 4
print(c)               # [1, 2]
\end{code}
\zam В список попадают \emph{значения} переменных на момент создания; изменение \texttt{b} потом
на список не влияет.

\lead{Длина списка:}
\begin{code}
a = [1, 2, 3]
print(len(a))          # 3
\end{code}

\lead{Удаление элемента по индексу:}
\begin{code}
del a[1]
print(a)               # [1, 3]
\end{code}

\lead{Соединение списков} (сцепление, сложение):
\begin{code}
a1 = [1, 2]
a2 = [3, 4]
print(a1 + a2)         # [1, 2, 3, 4]
lst = [1, 2, 3, 4]
lst += [99]
print(lst)             # [1, 2, 3, 4, 99]
\end{code}

\lead{Повтор элементов списка:}
\begin{code}
a = ['k'] * 3
print(a)               # ['k', 'k', 'k']
\end{code}

\lead{Принадлежность элемента списку} (\texttt{in}, \texttt{not in}):
\begin{code}
a = [1, 2, 44]
if 44 in a:
    print('содержится')
if 55 not in a:
    print('не содержится')
\end{code}
\vspace{-4pt}
{\small Во втором условии стоит \texttt{not in}\ispr{В рукописи \texttt{if 55 in a: print('не')}. Числа 55 в
списке нет, и сообщение не напечатается; судя по его тексту, имелось в виду \texttt{not in}.}.}

\subsection{Перебор элементов списка}
\lead{1 способ~--- через индексы:}
\begin{code}
for i in range(len(a)):
    print(a[i], i)
# 1 0
# 2 1
# 44 2
\end{code}
\lead{2 способ~--- через перебор элементов коллекции:}
\begin{code}
for el in a:
    print(el)
\end{code}

\subsection{Ввод списка с клавиатуры}
\lead{1 способ}~--- начинаем с пустого списка и добавляем элементы:
\begin{code}
a = []
n = int(input('Вв. количество элементов: '))
print(f'Вв. {n} эл-в')
for i in range(n):
    a.append(int(input(f'a[{i+1}] = ')))
print(a)
\end{code}
\lead{2 способ}~--- сразу создаём список нужной длины и заполняем по индексам:
\begin{code}
n = int(input())
print(n)
a = [0] * n            # заполнение массива a нулями
for i in range(n):
    a[i] = int(input())
print(a)
\end{code}

\subsection{Вывод элементов списка}
\begin{minipage}[t]{0.48\textwidth}
\begin{code}
print('Массив:')
for i in range(n):
    print(a[i], end=' ')
print()
\end{code}
\end{minipage}\hfill
\begin{minipage}[t]{0.48\textwidth}
\begin{code}
print('Массив:')
for el in a:
    print(el, end='\t')
print()
\end{code}
\end{minipage}

\texttt{end=' '}~--- после элемента печатается пробел, а не перевод строки; пустой \texttt{print()} в
конце переводит строку.

\subsection{Срезы списков}
\begin{center}
\texttt{<список>[<нач. инд.>:<инд. след. за кон.>:<шаг>]}
\end{center}
\begin{code}
a = [1, 2, 3, 4, 5]
print(a[0:len(a):2])   # [1, 3, 5]
print(a[::2])          # [1, 3, 5]
print(a[::-1])         # [5, 4, 3, 2, 1]
\end{code}
То же, что \texttt{a[::-1]}, циклом:
\begin{code}
b = []
for i in range(len(a) - 1, -1, -1):
    b.append(a[i])
print(b)               # [5, 4, 3, 2, 1]
\end{code}
Срезу можно присвоить новые значения:
\begin{code}
month = ['январь', 'февраль', 'март', 'апрель', 'май', 'июнь',
         'июль', 'август', 'сентябрь', 'октябрь', 'ноябрь', 'декабрь']
spring = month[2:5]
print(spring)          # ['март', 'апрель', 'май']
month[5:8] = ('June', 'July', 'August')
print(month[4:9])      # ['май', 'June', 'July', 'August', 'сентябрь']
\end{code}
\begin{code}
print(month)           # изменённый массив
a = [1, 2, 3, 4, 5, 6]
del a[::2]
print(a)               # [2, 4, 6]
print(a[1:99])         # [4, 6]   (выход за границы в срезе - не ошибка)
print(a[99:1000])      # []
\end{code}

\subsection{Генератор списков}
\begin{code}
a1 = [k for k in range(7)]
print(a1)              # [0, 1, 2, 3, 4, 5, 6]
a2 = [k ** 2 for k in range(5)]
print(a2)              # [0, 1, 4, 9, 16]
a3 = [k for k in range(7) if k % 3]
print(a3)              # [1, 2, 4, 5]
import random
a4 = [random.random() for _ in range(5)]
print(a4)              # 5 случайных чисел из [0; 1)
\end{code}
\zam В \texttt{a3} условие \texttt{k \% 3} истинно, когда остаток не равен нулю, т.\,е. берутся числа,
не делящиеся на~3. Имя \texttt{\_} ставят вместо параметра цикла, если он в теле не нужен.

\subsection{Функция \texttt{enumerate()}}
\texttt{enumerate()} даёт сразу и номер, и значение элемента:
\begin{code}
a1 = [2, 456, 88, 7]
for index, value in enumerate(a1, start=1):
    print(index, value)
# 1 2
# 2 456
# 3 88
# 4 7
\end{code}
Без \texttt{start} (\texttt{for index, value in enumerate(a1):}) нумерация идёт с нуля.

\subsection{Методы списков}
\begin{center}
\begin{tabular}{>{\ttfamily}l l}
append(x)    & добавляет элемент в конец списка\\
extend(lst)  & добавляет в конец все элементы другого списка\\
insert(i, a) & вставляет на $i$-ю позицию элемент $a$\\
remove(a)    & удаляет первое найденное значение элемента $a$; ошибка, если такого нет\\
pop(i)       & удаляет $i$-й элемент и возвращает его; без индекса~--- последний\\
index(a)     & возвращает индекс первого элемента, равного $a$\\
count(a)     & возвращает количество элементов, равных $a$\\
sort()       & сортирует список (меняет сам список)
\end{tabular}
\end{center}

\begin{code}
lst = [1, 2, 3, 4]
lst.insert(2, 'a')
print(lst)             # [1, 2, 'a', 3, 4]

a = [1, 2, 3, 1]
a.remove(1)
print(a)               # [2, 3, 1]

L = [1, 2, 3]
el = L.pop(1)
print(el, L)           # 2 [1, 3]

L = [1, 2, 3, 1]
i1 = L.index(1)
print(i1)              # 0
i2 = L.index(1, i1 + 1)    # искать, начиная с позиции i1 + 1
print(i2)              # 3
\end{code}

\begin{code}
L = [1, 2, 3, 4, 1, 6, 1]
print(L.count(1))      # 3
# то же «вручную» - сколько раз встречается 4:
k = 0
for el in L:
    if el == 4:
        k += 1
print(k)               # 1
\end{code}

\begin{code}
L = [1, 56, -8, 1000, 2]
L.sort()
print(L)               # [-8, 1, 2, 56, 1000]
# функция sorted() не меняет исходный список
L = [4, 7, 1]
print(sorted(L, reverse=True))     # [7, 4, 1]
print(L)                           # [4, 7, 1]
L = [-10, 1, -5, 2, -3]
print(sorted(L, key=abs))          # [1, 2, -3, -5, -10]
\end{code}
\texttt{reverse=True}~--- по убыванию; \ \texttt{key=abs}~--- сравнивать элементы по модулю.

\begin{code}
L = [1, 2, 3]
L.reverse()            # reverse() - разворачивает список
print(L)               # [3, 2, 1]
print(L[::-1])         # [1, 2, 3]
\end{code}

\lead{Метод \texttt{copy()}}~--- поверхностная копия списка:
\begin{code}
L1 = [1, 2, 3]
L2 = L1.copy()         # сравни: L2 = L1
L1.reverse()
print(L2, id(L1), id(L2))     # [1, 2, 3], id(L1) != id(L2)
\end{code}
\vspace{-4pt}
{\small Разворачивается список \texttt{L1}\ispr{В рукописи \texttt{L1 = reverse()}~--- такой функции нет,
нужен метод \texttt{L1.reverse()}.}.}

При \texttt{L2 = L1} будет \texttt{id(L1) == id(L2)}. Поэтому при использовании \texttt{copy()} значение
нового списка не будет изменяться при изменении исходного ($id_1\ne id_2$). Если списки присвоить
через \texttt{=} (присваивание), то значение копии будет изменяться при изменении оригинала
($id_1=id_2$): обе переменные указывают на один и тот же список. Копию можно сделать и срезом:
\texttt{L3 = L1[:]}.

\lead{Метод \texttt{clear()}}~--- очистка списка:
\begin{code}
print(id(L1))          # id1
L1.clear()
print(L1, id(L1))      # [] id2:  id1 == id2 - тот же список
L1 = []
print(L1, id(L1))      # [] id3:  id1 != id3 - новый список
del L1[:]              # тоже очищает список на месте
print(id(L1))
\end{code}

\primn{9.1} Нахождение среднего значения элементов одномерного массива.
\begin{code}
# Переменные:
#     i - индекс элемента
#     n - размер массива
#     a - одномерный массив
#     S - сумма
#     Sr - среднее
n = int(input('Введите размер мас.: '))

a = []
print(f'Введите {n} эл-в массива')
for i in range(n):
    print(f'a[{i+1}] = ', end='')
    a.append(float(input()))
# вывод массива
print('массив:')
for el in a:
    print(el, end=' ')
print()

# определение суммы элементов массива
S = 0
for el in a:
    S += el
Sr = S / n
print(f'Сумма элементов массива {S:.2f}')
print(f'Среднее значение {Sr:.2f}')
\end{code}

\lead{Ввод элементов массива одной строкой} (через пробел):
\begin{code}
lst = list(map(float, input().split()))
print(sum(lst) / len(lst))
\end{code}
\texttt{input().split()} разбивает строку на части по пробелам, \texttt{map(float, \dots)} переводит
каждую часть в число, \texttt{list(\dots)} собирает результат в список.

\begin{schema}
Схема к примеру 9.1. Справа от цикла~--- его краткая запись одним блоком; так и рисуют, когда
ввод или вывод массива ничем не примечателен.
\begin{center}
\begin{tikzpicture}[sa]
\node[term] (b) {начало};
\node[proc,below=of b] (p0) {$a$~--- пустой массив};
\node[io,below=of p0] (in) {Ввод $n$};
\node[prep,below=6mm of in] (h1) {$i=0(1)n-1$};
\node[io,below=of h1] (o1) {Вывод $i+1$};
\node[io,below=of o1] (i1) {Ввод и добавление\\$a_i$ в конец};
\draw (b)--(p0)--(in)--(h1)--(o1)--(i1);
\cycleold{h1}{i1}
\node[prep,below=6mm of h1-out.center] (h2) {$el$ из $a$};
\node[io,below=of h2] (o2) {Вывод $el$};
\draw (h1-out.center)--(h2)--(o2);
\cycleold{h2}{o2}
\node[draw,circle,minimum size=7mm,inner sep=0pt,below=6mm of h2-out.center] (c1) {1};
\draw[->] (h2-out.center)--(c1);
% краткая запись
\draw[decorate,decoration={brace,amplitude=6pt}] ([xshift=2.5cm]h1.north) -- ([xshift=2.5cm]h1-out.center);
\node[io,minimum width=2.8cm,anchor=west] at ([xshift=2.9cm]$(h1.north)!0.5!(h1-out.center)$) {Ввод\\массива $a$};
\draw[decorate,decoration={brace,amplitude=6pt}] ([xshift=2.5cm]h2.north) -- ([xshift=2.5cm]h2-out.center);
\node[io,minimum width=2.8cm,anchor=west] at ([xshift=2.9cm]$(h2.north)!0.5!(h2-out.center)$) {Вывод\\массива $a$};
% продолжение
\begin{scope}[xshift=-6.2cm,yshift=-1.0cm]
\node[draw,circle,minimum size=7mm,inner sep=0pt] (c2) {1};
\node[proc,below=6mm of c2] (s0) {$S=0$};
\node[prep,below=6mm of s0] (h3) {$el$ из $a$};
\node[proc,below=of h3] (s1) {$S=S+el$};
\draw[->] (c2)--(s0); \draw (s0)--(h3)--(s1);
\cycleold{h3}{s1}
\node[proc,below=6mm of h3-out.center] (sr) {$S_r=\dfrac{S}{n}$};
\node[io,below=of sr] (out) {Вывод $S,\ S_r$};
\node[term,below=of out] (e) {конец};
\draw (h3-out.center)--(sr)--(out)--(e);
\end{scope}
\end{tikzpicture}
\end{center}
\end{schema}

\primn{9.2} Сортировка массива по убыванию методом пузырька.
\begin{code}
# Переменные:
#     a - одномерный массив (список)
#     n - размер массива
#     i, k - параметры циклов
#     temp - вспомогательная переменная для обмена элементов
import random

# заполнение массива случайными числами
n = int(input('Вв. кол-во эл-в массива: '))
a = [0] * n
for i in range(n):
    a[i] = random.randint(-10, 10)

# вывод массива
for i in range(n):
    print(a[i], end='\t')
print()

# сортировка пузырьком
for i in range(n - 1):
    for k in range(n - i - 1):
        if a[k] < a[k + 1]:          # > - по возрастанию
            temp = a[k]
            a[k] = a[k + 1]
            a[k + 1] = temp
            # или: a[k], a[k + 1] = a[k + 1], a[k]
print('Отсортированный массив:')
print(*a, sep='\t')
\end{code}
\zam За один проход внутреннего цикла самый маленький из ещё не упорядоченных элементов
<<всплывает>> в конец, поэтому каждый следующий проход на один элемент короче:
$k=0(1)n-i-2$. Звёздочка в \texttt{print(*a, sep='\textbackslash t')} передаёт элементы списка в
\texttt{print} по отдельности.

\begin{center}
\begin{tikzpicture}[sa]
\node[term] (b) {начало};
\node[io,below=of b] (in) {Ввод $n$};
\node[proc,below=of in] (z) {Заполнение массива\\$a$ $n$ нулями};
\node[prep,below=6mm of z] (h1) {$i=0(1)n-1$};
\node[proc,below=of h1] (r) {$a_i$~--- случайное\\число от $-10$ до $10$};
\draw (b)--(in)--(z)--(h1)--(r);
\cycleold{h1}{r}
\node[io,below=6mm of h1-out.center] (o1) {Вывод\\массива $a$};
\draw (h1-out.center)--(o1);
\node[draw,circle,minimum size=7mm,inner sep=0pt,below=6mm of o1] (c1) {1};
\draw[->] (o1)--(c1);
\draw[decorate,decoration={brace,amplitude=6pt}] ([xshift=2.5cm]h1.north) -- ([xshift=2.5cm]h1-out.center);
\node[io,minimum width=2.8cm,anchor=west] at ([xshift=2.9cm]$(h1.north)!0.5!(h1-out.center)$) {Ввод\\массива $a$};
\end{tikzpicture}
\end{center}
\begin{center}
\begin{tikzpicture}[sa]
\node[draw,circle,minimum size=7mm,inner sep=0pt] (c2) {1};
\node[prep,below=6mm of c2] (hi) {$i=0(1)n-2$};
\node[prep,below=6mm of hi] (hk) {$k=0(1)n-i-2$};
\node[dec,below=6mm of hk] (d) {$a_k<a_{k+1}$};
\node[proc,minimum width=2.6cm] (sw) at ([xshift=2.4cm,yshift=-1.5cm]d.center)
  {$temp=a_k$\\$a_k=a_{k+1}$\\$a_{k+1}=temp$};
\draw[->] (c2)--(hi); \draw (hi)--(hk)--(d);
\draw[->] (d.east) -| node[yes,above,pos=0.25]{да} (sw.north);
\coordinate (m) at ($(d.center |- sw.south)+(0,-0.35)$);
\draw (d.west) -- node[yes,above]{нет} ++(-0.4,0) |- (m);
\draw (sw.south) |- (m);
\cycleold[2.5]{hk}{m}
\cycleold[2.9]{hi}{hk-out}
\node[io,below=6mm of hi-out.center] (o2) {Вывод\\массива $a$};
\node[term,below=of o2] (e) {конец};
\draw (hi-out.center)--(o2)--(e);
\end{tikzpicture}
\end{center}

\primn{9.3} Поиск максимума и минимума в массиве.
\begin{code}
# Переменные:
#     a - массив
#     n - размер массива
#     el - элемент массива
#     i - индекс элемента массива
#     amax - макс. элемент
#     amin - мин. элемент
#     imin - его индекс
a = list(map(int, input('Введите элементы мас.: ').split()))
print(*a)

amax = -float('inf')
for el in a:
    if el > amax:
        amax = el
print('Макс. элемент', amax)

# поиск минимума и его индекса
amin = a[0]
imin = 0
n = len(a)
for i in range(1, n):
    if a[i] < amin:
        amin = a[i]
        imin = i
print(f'min элемент {amin}, его номер {imin + 1}')
\end{code}
\zam Максимум можно начинать с $-\infty$ (\texttt{-float('inf')}): любой элемент больше, и первое же
сравнение заменит его. Минимум здесь начинают с первого элемента, тогда цикл идёт с $i=1$.
Номер для человека~--- \texttt{imin + 1}, т.\,к. индексы считаются с нуля.

% =====================================================================
\section{Матрицы (двумерные массивы)}
% =====================================================================
\begin{hwbox}
В рукописи матрицы записаны после кортежей (тема~11), здесь они стоят по номеру темы. У заголовка
пометка <<без п/п>>: примеры этой темы написаны без подпрограмм.
\end{hwbox}

\begin{opr}
\textbf{Определение.} Если число индексов в описании массива равно~2, то массив называется
\emph{двумерным массивом} или \emph{матрицей}.
\end{opr}
Матрицу в Питоне можно представить в виде вложенного списка, при этом для обращения к элементу
нужно указывать 2~индекса, каждый в своих скобках \texttt{[ ]}.

\begin{center}
\begin{tikzpicture}[line width=0.6pt,>={Stealth[length=2mm]}]
\node (A) at (0,0) {$A=\begin{pmatrix}11&12&13&14\\21&22&23&24\\31&32&33&34\end{pmatrix}$};
\draw[->] (2.35,0.75) -- node[right,align=left,font=\small]{строки~--- индекс $i$:\\
$i=0,1,\dots,n-1$} (2.35,-0.75);
\draw[->] (-0.45,-1.0) -- node[below,font=\small]{столбцы~--- индекс $j$: \ $j=0,1,\dots,m-1$} (2.0,-1.0);
\end{tikzpicture}
\end{center}
$A_{n\times m}$: \ $n$~--- число строк, \ $m$~--- число столбцов.

\begin{code}
A = [[11, 12, 13, 14],
     [21, 22, 23, 24],
     [31, 32, 33, 34]]
# A[i][j] - элемент из i-й строки и j-го столбца
print(A[1][2])         # 23
\end{code}

\subsection{Создание и ввод матрицы}
Создание матрицы $A_{n\times m}$ через генератор и заполнение её нулями:
\begin{code}
A = [[0] * m for _ in range(n)]
\end{code}
\imp Не пиши \texttt{A = [[0] * m] * n}: так получится $n$ ссылок на \emph{одну и ту же} строку, и
присваивание \texttt{A[0][0] = 5} изменит первый элемент во всех строках.

\lead{Ввод элементов матрицы:}
\begin{code}
n = int(input('Вв. количество строк: '))
m = int(input('Вв. количество столбцов: '))
A = [[0] * m for _ in range(n)]
print(f'Введите {n * m} элементов матрицы')
for i in range(n):
    for j in range(m):
        A[i][j] = float(input())
print(A)
\end{code}

\lead{Ввод по строкам:} каждая строка матрицы вводится одной строкой через пробел
\begin{code}
B = []
for i in range(n):
    row = [int(el) for el in input().split()]
    B.append(row)
print(B)
\end{code}
или по одному элементу:
\begin{code}
B = []
for i in range(n):
    row = []
    for j in range(m):
        row.append(int(input()))
    B.append(row)
print(B)
\end{code}

\subsection{Вывод матрицы}
\begin{minipage}[t]{0.48\textwidth}
\begin{code}
for i in range(n):
    for j in range(m):
        print(A[i][j], end='\t')
    print()
\end{code}
\end{minipage}\hfill
\begin{minipage}[t]{0.48\textwidth}
\begin{code}
for row in A:
    for el in row:
        print(el, end='\t')
    print()
\end{code}
\end{minipage}

\subsection{Пример}
\primn{10.1} Найти в матрице максимальный элемент, минимальный элемент и его индексы. Заменить
минимальный элемент максимальным.
\begin{code}
# Переменные:
#     a - матрица
#     i, j - параметры циклов
#     n, m - число строк и столбцов матрицы
#     amin - мин. элемент
#     imin, jmin - его индексы
#     amax - макс. элемент

# ввод матрицы (см. п. 10.1)
# вывод матрицы (см. п. 10.2)

# поиск минимума и максимума
amin = a[0][0]
imin = 0
jmin = 0
amax = -float('inf')
for i in range(n):
    for j in range(m):
        if a[i][j] > amax:
            amax = a[i][j]
        if a[i][j] < amin:
            amin = a[i][j]
            imin = i
            jmin = j
print(f'Макс. элемент {amax}')
print(f'Мин. элемент {amin}, его индексы [{imin}][{jmin}]')

# замена минимума на максимум
a[imin][jmin] = amax

# вывод матрицы
print('Матрица после замены:')
for row in a:
    for el in row:
        print(el, end='\t')
    print()
\end{code}

\begin{schema}
На схеме двойной цикл ввода (вывода) матрицы можно заменить одним блоком:
\begin{center}
\begin{tikzpicture}[sa]
\node[prep] (hi) at (0,0) {$i=0(1)n-1$};
\node[prep,below=6mm of hi] (hj) {$j=0(1)m-1$};
\node[io,below=of hj] (in) {Ввод $a_{ij}$};
\draw (0,0.9)--(hi); \draw (hi)--(hj)--(in);
\cycleold{hj}{in}
\cycleold[0.9]{hi}{hj-out}
\draw (hi-out.center) -- ++(0,-0.4);
\draw[decorate,decoration={brace,amplitude=6pt}] (2.6,0.5) -- (2.6,-4.6);
\node[io,minimum width=2.4cm] at (4.25,-2.05) {Ввод\\матрицы $a$};
\begin{scope}[xshift=8.9cm]
\node[prep] (oi) at (0,0) {$i=0(1)n-1$};
\node[prep,below=6mm of oi] (oj) {$j=0(1)m-1$};
\node[io,below=of oj] (out) {Вывод $a_{ij}$};
\draw (0,0.9)--(oi); \draw (oi)--(oj)--(out);
\cycleold{oj}{out}
\cycleold[0.9]{oi}{oj-out}
\draw (oi-out.center) -- ++(0,-0.4);
\draw[decorate,decoration={brace,amplitude=6pt}] (2.6,0.5) -- (2.6,-4.6);
\node[io,minimum width=2.4cm] at (4.25,-2.05) {Вывод\\матрицы $a$};
\end{scope}
\end{tikzpicture}
\end{center}
\end{schema}

\begin{center}
Схема к примеру 10.1\dop{В рукописи схема не закончена: после блока начальных значений стоят пустые
блоки циклов и решения. Здесь она достроена по программе.}\\*[4pt]
\begin{tikzpicture}[sa,bx/.style={proc,minimum width=2.6cm}]
% левый столбец
\node[term] (b) at (0,0) {начало};
\node[io,below=of b] (in) {Ввод $n,\ m$};
\node[proc,below=of in] (z) {Заполнение матрицы\\$a_{n\times m}$ нулями};
\node[io,below=of z] (ia) {Ввод\\матрицы $a$};
\node[io,below=of ia] (oa) {Вывод\\матрицы $a$};
\node[proc,below=of oa] (st) {$a_{\min}=a_{00}$\\$i_{\min}=0$, \ $j_{\min}=0$\\$a_{\max}=-\infty$};
\node[draw,circle,minimum size=7mm,inner sep=0pt,below=6mm of st] (c1) {1};
\draw (b)--(in)--(z)--(ia)--(oa)--(st); \draw[->] (st)--(c1);
% правый столбец
\begin{scope}[xshift=7.2cm]
\node[draw,circle,minimum size=7mm,inner sep=0pt] (c2) at (0,0) {1};
\node[prep,below=6mm of c2] (hi) {$i=0(1)n-1$};
\node[prep,below=6mm of hi] (hj) {$j=0(1)m-1$};
\node[dec,below=6mm of hj] (d1) {$a_{ij}>a_{\max}$};
\node[bx] (s1) at ([xshift=2.6cm,yshift=-1.3cm]d1.center) {$a_{\max}=a_{ij}$};
\coordinate (m1) at ($(d1.center |- s1.south)+(0,-0.35)$);
\draw[->] (d1.east) -| node[yes,above,pos=0.25]{да} (s1.north);
\draw (d1.west) -- node[yes,above]{нет} ++(-0.4,0) |- (m1);
\draw (s1.south) |- (m1);
\node[dec,below=5mm of m1] (d2) {$a_{ij}<a_{\min}$};
\node[bx] (s2) at ([xshift=2.6cm,yshift=-1.55cm]d2.center) {$a_{\min}=a_{ij}$\\$i_{\min}=i$\\$j_{\min}=j$};
\coordinate (m2) at ($(d2.center |- s2.south)+(0,-0.35)$);
\draw[->] (d2.east) -| node[yes,above,pos=0.25]{да} (s2.north);
\draw (d2.west) -- node[yes,above]{нет} ++(-0.4,0) |- (m2);
\draw (s2.south) |- (m2);
\draw (m1)--(d2);
\draw[->] (c2)--(hi); \draw (hi)--(hj)--(d1);
\cycleold[2.75]{hj}{m2}
\cycleold[3.15]{hi}{hj-out}
\node[io,below=6mm of hi-out] (o1) {Вывод $a_{\max}$,\\$a_{\min},\ i_{\min},\ j_{\min}$};
\node[proc,below=of o1] (r) {$a_{i_{\min}j_{\min}}=a_{\max}$};
\node[io,below=of r] (o2) {Вывод\\матрицы $a$};
\node[term,below=of o2] (e) {конец};
\draw (hi-out.center)--(o1)--(r)--(o2)--(e);
\end{scope}
\end{tikzpicture}
\end{center}

% =====================================================================
\section{Кортежи}
% =====================================================================
\begin{opr}
\textbf{Определение.} \emph{Кортеж}~--- это упорядоченная неизменяемая коллекция, используемая для
хранения набора элементов. \emph{Упорядоченная} означает, что у каждого элемента есть свой
порядковый номер (индекс).
\end{opr}
Отличие кортежа от списка состоит в том, что кортежи нельзя изменять.

Ключевое слово, отвечающее за данную коллекцию,~--- \texttt{tuple}. Элементы кортежа заключаются в
круглые скобки, типы данных могут быть любыми.
\begin{code}
card = ('7', 'пики')
empty1 = ()            # пустой кортеж
empty2 = tuple()
\end{code}

\imp Для создания кортежа из одного элемента необходимо поставить после этого элемента запятую:
\begin{code}
t1 = (18,)
t2 = ('one',)
\end{code}

Если не возникает неоднозначности, элементы кортежа можно перечислить без скобок:
\begin{code}
k1 = (13, 22, 444, -5)
k2 = 13, 22, 444, -5
print(k1 == k2)        # True
\end{code}
\zam В рукописи записано \texttt{id(k1) = id(k2)}. В программе так и выходит: Питон хранит два
одинаковых кортежа-константы как один объект. Но это особенность интерпретатора, а не правило,
поэтому кортежи сравнивают через~\texttt{==}.

\begin{code}
print(card[0], len(card), card + t1)
# 7 2 ('7', 'пики', 18)
\end{code}
Для кортежей тоже есть операция среза, она возвращает новый кортеж.

К кортежам применимы операции сравнения, сравнение производится поэлементно:
\begin{code}
(1, 2) == (1, 3)                  # False
(1, 2) < (1, 3)                   # True
(1, 2) < (5,)                     # True
('7', 'червей') < ('7', 'треф')   # False
(1, 2) < ('7', 'пик')             # ошибка: число нельзя сравнить со строкой
\end{code}
Для обработки кортежей можно использовать функции и методы списков, кроме тех, которые изменяют
элементы.

\subsection{Присваивание кортежей}
Кортежи можно присваивать друг другу:
\begin{code}
n, s = 10, 'hello'
a, b = b, a            # обмен значений
\end{code}

\subsection{Для чего нужны кортежи}
\begin{enumerate}[label=\arabic*)]
\item Набор данных (список), который нужно защитить от изменений.
\item Обмен значений переменных местами.
\item Если в цикле есть набор переменных, которые изменяются одновременно.
\item Кортежи занимают меньше места в памяти.
\item Кортежи работают быстрее.
\item В функции~--- чтобы вернуть несколько значений:
\begin{code}
def f(a, b):
    return a + b, a - b

y1, y2 = f(5, 2)
k = f(5, 2)
print(k[0], k[1])      # 7 3
\end{code}
\item Для передачи переменного числа аргументов в п/п можно использовать кортеж:
\begin{code}
def proc(*args):
    print(type(args))  # класс tuple
    for value in args:
        print(value)

proc(1, 2, 3)
\end{code}
\item Кортеж может быть использован в словаре в качестве составного ключа, например координат
или ФИО.
\item Кортеж может быть параметром цикла \texttt{for}:
\begin{code}
lst = [11, 2, 43, 4]
for i, el in enumerate(lst, start=1):
    print(i, el)
\end{code}
или пример для пункта~8\dop{Словарь \texttt{locations} добавлен, чтобы пример можно было
запустить.}:
\begin{code}
locations = {(55.76, 37.62): 'Москва', (54.63, 39.74): 'Рязань'}
for coord, city in locations.items():
    print(f'{city} на координатах {coord}')
\end{code}
\item В некоторых библиотеках параметры п/п тоже могут быть кортежами:
\begin{code}
import turtle
turtle.write('Привет!', font=('Arial', 16, 'bold'))
\end{code}
\item Для ввода нескольких переменных одной строкой:
\begin{code}
a, b = map(int, input().split())
\end{code}
\end{enumerate}

\subsection{Ввод кортежа с клавиатуры}
\begin{code}
n = int(input('Введите число элементов кортежа: '))
k = tuple(int(input(f'Введите число {i + 1}: ')) for i in range(n))
print(k)
\end{code}

\end{document}
